\section{The subcritical depth threshold}

The coefficient family has \(O(L^2)\) members, but its total parameter length
is only linear in \(L\).
\begin{lemma}
For \(M=2L-1\),
\[
 \sum_{1\le a,b\le M}\frac1{\max(a,b)}
 =\sum_{m\le M}\frac{2m-1}{m}<4L. \tag{9}
\]
\end{lemma}
\begin{proof}
Exactly \(2m-1\) ordered pairs have maximum \(m\).  Summing their weights
gives \(2M-\sum_{m\le M}1/m<2M<4L\).
\end{proof}

\begin{theorem}[Growing-depth incidence bound]
For fixed \(A>0\), uniformly over \(1\le L\le(\log Q)^A\),
\[
 C_L(Q)\ll_A
 \frac{LQ\log\log(3LQ)}{(\log Q)^2}. \tag{10}
\]
\end{theorem}
\begin{proof}
Sum (8) over \(a,b<2L\) and apply (9).  The summed constant and remainder
terms are \(o_A(LQ/\log^2Q)\), because \(L\) is polylogarithmic.
\end{proof}

\begin{corollary}[Subcritical obstruction]
If \(L=o(\log Q/\log\log Q)\), then
\[
 \frac{C_L(Q)}{|\PP_Q|}\longrightarrow0. \tag{11}
\]
If positive row masses have a fixed maximum-to-minimum ratio \(\kappa\), the
maximum saving supported only on these collisions is also \(o(1)\).
\end{corollary}
\begin{proof}
The prime number theorem gives \(|\PP_Q|\sim Q/\log Q\), which proves (11).
At most \(2C_L(Q)\) rows are incident to a collision, so their mass fraction
is at most \(2\kappa C_L(Q)/|\PP_Q|\).  All nonincident rows remain
orthogonal to every other row and retain their diagonal energy.  Hence the
possible saving is bounded by the incident mass fraction.
\end{proof}

In particular, no fixed \(\delta>0\), including \(\delta=1/400\), can be paid
at subcritical depth.  The corollary is only a necessary condition:
\(L\asymp\log Q/\log\log Q\) is not proved sufficient.
